% generated from papers/VI-nonabelian/paper_VI_nonabelian.tex -- edit the paper, then rerun assemble.py
\chapter[The non-abelian obstruction]{Towards non-abelian localized Bost--Connes systems:\\ the matrix Kakutani dichotomy and a Tannakian obstruction}\label{ch:VI}
\chapterorigin{This chapter is Paper VI of the series (\texttt{papers/VI-nonabelian/}).}
\begin{summary}
The $GL_1$ theory of [Ch.~\ref{ch:Main}] classifies the $\mathrm{KMS}_\beta$ simplex of
$\cA_{K,S}$ as $\Prob(G_S/\Xi_\beta^\perp)$, where $\Xi_\beta$ is a subgroup of the
Pontryagin dual $\widehat{G_S}$ cut out by a Dirichlet series condition. We examine what
survives when the abelian symmetry group is replaced by a Galois group with non-abelian
image, and separate three things that are usually conflated.

\emph{What transfers.} The Dirichlet-series half of the Kakutani dichotomy is
representation-theoretic, not character-theoretic. For a unitary representation $\rho$ of a
compact group equipped with Frobenius classes we take as the obstruction quantity the
Hilbert--Schmidt Frobenius series
\[
\Xi_\beta(\rho,S)=\sum_{\pp\in S}\Ns\pp^{-\beta}\bigl\|I_n-\rho(\Frob_\pp)\bigr\|_{\HS}^2
=2\sum_{\pp\in S}\Ns\pp^{-\beta}\bigl(n-\Re\Tr\rho(\Frob_\pp)\bigr),
\]
which reduces to the abelian condition for characters, and whose convergence is the
coboundary condition for the $U(n)$-valued Frobenius cocycle whenever that cocycle is
defined, namely when the Frobenius classes commute; without commutation there is no
cocycle on the tail relation at all, and the series is the primary object. Its relation
to $L$-functions is
\[
\Xi_\beta(\rho,S)=2n\log\zeta_{K,S}(\beta)-2\Re\log L_S(\beta,\rho)+O(1),
\]
so that at $\beta\le1$ and $S=P_K$ the obstruction is absent for every nontrivial
irreducible $\rho$, the input being $L(1,\rho)\ne0,\infty$: Hecke--Landau for Artin
representations, Jacquet--Shalika in the automorphic case. For Artin representations the
Chebotarev criterion of [Ch.~\ref{ch:Main}, Thm.~3.9] transfers with
$D_\rho=\{\pp:\rho(\Frob_\pp)\ne I\}$.

\emph{What does not transfer, and its replacement.} The classification
$\Prob(G_S/\Xi_\beta^\perp)$ uses Pontryagin duality twice and fails without it: for
non-abelian $G$ the set of irreducible representations is not a group, $\Xi_\beta$ cannot be
a subgroup of it, and $\Xi_\beta^\perp$ is undefined. The correct replacement is not a
subgroup of a dual but a \emph{quotient of the group}. We prove that
$\cT_\beta(S)=\{\rho:\Xi_\beta(\rho,S)<\infty\}$ is closed under direct sums, subobjects,
tensor products and duals, hence is a Tannakian subcategory of $\Rep(G)$, and therefore
equals $\Rep(G/N_\beta)$ for a closed normal subgroup $N_\beta=\bigcap_{\rho\in\cT_\beta}\ker\rho$.
In the abelian case $N_\beta=\Xi_\beta^\perp$ and the two formulations agree.

\emph{What is not proved here.} We do not construct a $GL_n$ algebra whose
$\mathrm{KMS}_\beta$ simplex is $\Prob(G/N_\beta)$, and we explain precisely which step of
the $GL_1$ proof has no available substitute. Two concrete findings bear on the $GL_2$ case,
formulated for the Sato--Tate classes of a non-CM elliptic curve. If the classes along $S$
are Sato--Tate equidistributed then $\Xi_\beta(\rho,S)\sim4\sum_{\pp\in S}\Ns\pp^{-\beta}$
diverges whenever the partition function does, so \emph{no} symmetry breaking arises from a
Sato--Tate representation. Conversely, assuming Sato--Tate as an equidistribution input, we
construct sparse $S$ with $\zeta_S(1)=\infty$ and $\Xi_1(\rho,S)<\infty$ for the
representation attached to a non-CM elliptic curve, by selecting blocks of primes with
$\theta_p\le2^{-j}$; the exotic phase is therefore available, but only on deliberately
constructed anomalous sets, never generically.
\end{summary}

\section{Which steps of the abelian argument use commutativity}

Recall the $GL_1$ chain of reasoning [Ch.~\ref{ch:Main}]. A $\mathrm{KMS}_\beta$ state is a measure
$\nu$ on $Y_S$ scaling by $\Ns\aaa^{-\beta}$; the symmetry group $G_S$ acts on the simplex;
a character $\chi\in\widehat{G_S}$ fails to be an obstruction exactly when $\chi\circ\sigma$
is a coboundary for $\pi_\beta$, which by Kakutani is the convergence of
$\sum_\pp\Ns\pp^{-\beta}|1-\chi(\sigma_\pp)|^2$; the obstruction characters form a
\emph{subgroup} $\Xi_\beta\le\widehat{G_S}$; and the simplex is
$\Prob(G_S/\Xi_\beta^\perp)$.

Three of these steps are insensitive to commutativity and two are not.

\begin{itemize}
\item[(i)] The Kakutani step has two halves: a Dirichlet series in Frobenius data, and its
identification with a coboundary condition for a cocycle on the tail relation.
\S\ref{VI:sec:kakutani} shows that the series transfers, with the Hilbert--Schmidt norm in
place of $|1-z|$, and that the cocycle half transfers exactly when the Frobenius classes
along $S$ commute --- the tail relation being generated by commuting shifts, a
$U(n)$-valued cocycle with prescribed non-commuting increments does not exist.
\item[(ii)] The passage from convergence to a group structure uses that $\widehat{G_S}$ is a
group under pointwise product. For non-abelian $G$ the irreducibles are not closed under
tensor product as a group law: $\rho_1\otimes\rho_2$ decomposes. \S\ref{VI:sec:tannaka} replaces
the subgroup by a Tannakian subcategory.
\item[(iii)] The identification of the simplex as $\Prob(G_S/\Xi_\beta^\perp)$ uses
Pontryagin duality to convert a subgroup of the dual into a closed subgroup of the group.
This is the step with no substitute at the level of duals; \S\ref{VI:sec:tannaka} supplies the
substitute at the level of the group itself.
\item[(iv)] The $L$-function dictionary is an Euler-product identity and transfers
(\S\ref{VI:sec:L}), with different constants from those one might guess.
\item[(v)] The construction of an algebra whose simplex realizes the answer does \emph{not}
transfer, and \S\ref{VI:sec:algebra} says why.
\end{itemize}

Throughout, $K$ is a number field, $S$ a set of finite primes, and $G$ a compact group
equipped with a family of conjugacy classes $\Frob_\pp\in G$ indexed by all but finitely
many $\pp\in S$: typically $G$ is a profinite quotient of $\Gal(\bar K/K)$ unramified outside
a finite set and $\Frob_\pp$ the Frobenius class, but $G$ may also be a Sato--Tate group and
$\Frob_\pp$ the Satake class, as in \S5. All representations are continuous, unitary and
finite-dimensional; sums over $\pp\in S$ omit the finitely many primes at which
$\Frob_\pp$ is undefined, which affects no convergence question. We write
$\zeta_{K,S}(\beta)=\prod_{\pp\in S}(1-\Ns\pp^{-\beta})^{-1}$ as in [Ch.~\ref{ch:Main}], so that
$\log\zeta_{K,S}(\beta)=\sum_{\pp\in S}\Ns\pp^{-\beta}+O(1)$ for $\beta>1/2$.

\section{The matrix Kakutani dichotomy}\label{VI:sec:kakutani}

\begin{definition}\label{VI:def:Xi}
For a representation $\rho:G\to U(n)$ put
\begin{equation}\label{VI:eq:Xi}
\Xi_\beta(\rho,S):=\sum_{\pp\in S}\Ns\pp^{-\beta}\bigl\|I_n-\rho(\Frob_\pp)\bigr\|_{\HS}^2 .
\end{equation}
\end{definition}

\begin{lemma}[Trace form]\label{VI:lem:trace}
For $U\in U(n)$, $\|I_n-U\|_{\HS}^2=2\bigl(n-\Re\Tr U\bigr)$. Hence
\begin{equation}\label{VI:eq:trace}
\Xi_\beta(\rho,S)=2\sum_{\pp\in S}\Ns\pp^{-\beta}\bigl(n-\Re\Tr\rho(\Frob_\pp)\bigr),
\end{equation}
a sum of non-negative terms, each in $[0,4n]$, depending only on the conjugacy class of
$\rho(\Frob_\pp)$.
\end{lemma}

\begin{proof}
$\|I-U\|_{\HS}^2=\Tr\bigl((I-U)^*(I-U)\bigr)=\Tr(I-U-U^*+U^*U)=\Tr(2I-U-U^*)=2n-2\Re\Tr U$,
using $U^*U=I$. Non-negativity and the bound follow from $|\Tr U|\le n$; conjugation
invariance of the trace gives the last claim.
\end{proof}

\begin{theorem}[Matrix Kakutani]\label{VI:thm:kakutani}
Let $\pi_\beta=\bigotimes_{\pp\in S}\mu_{t_\pp}$ with $t_\pp=\Ns\pp^{-\beta}$ as in
[Ch.~\ref{ch:Main}, Lem.~2.10], on $V=\N^{(S)}$ with the tail relation $\cR_S$, and put
$U_\pp=\rho(\Frob_\pp)\in U(n)$ (a representative of the class).
\begin{enumerate}
\item A $U(n)$-valued cocycle $c_\rho$ on $\cR_S$ with $c_\rho(v+e_\pp,v)=U_\pp$ for all
$\pp$ and $v$ exists if and only if the $U_\pp$, $\pp\in S$, commute pairwise.
\item Suppose they do. Then there is a measurable $g:V\to U(n)$ with
$g(v+e_\pp)=U_\pp\,g(v)$ a.e.\ for every $\pp$ --- i.e.\ $c_\rho$ is a coboundary --- if
and only if $\Xi_\beta(\rho,S)<\infty$.
\end{enumerate}
\end{theorem}

\begin{proof}
(1) The cocycle identity along the two paths from $v$ to $v+e_\pp+e_\qq$ gives
$c_\rho(v+e_\pp+e_\qq,v)=U_\qq U_\pp=U_\pp U_\qq$; conversely, for commuting $U_\pp$ the
formula $c_\rho(v+a,v)=\prod_\pp U_\pp^{a_\pp}$ is well defined and is a cocycle.

(2) Commuting unitaries are simultaneously diagonalizable: some unitary $W$ satisfies
\[
WU_\pp W^*=\mathrm{diag}\bigl(\chi_1(\pp),\dots,\chi_n(\pp)\bigr)\quad\text{for every }\pp,
\qquad\chi_r(\pp)\in\T .
\] Since $\|I-U_\pp\|_{\HS}^2=\sum_r|1-\chi_r(\pp)|^2$,
\[
\Xi_\beta(\rho,S)=\sum_{r=1}^n\sum_{\pp\in S}\Ns\pp^{-\beta}|1-\chi_r(\pp)|^2 .
\]
If $g$ exists, put $h=Wg$, so that $h(v+e_\pp)=\mathrm{diag}(\chi_r(\pp))\,h(v)$ and each
entry satisfies $h_{rs}(v+e_\pp)=\chi_r(\pp)h_{rs}(v)$. For each $r$ some entry $h_{rs}$ is
nonzero on a set of positive measure, $h(v)$ being unitary; by [Ch.~\ref{ch:Main}, Lem.~3.1]
$|h_{rs}|$ is a.e.\ constant, hence a.e.\ nonzero, and [Ch.~\ref{ch:Main}, Prop.~3.3], whose proof
uses only $|\chi_r(\pp)|=1$, gives the convergence of the $r$-th series. Conversely,
if all $n$ series converge, [Ch.~\ref{ch:Main}, Prop.~3.4] gives measurable $g_r:V\to\T$ with
$g_r(v+e_\pp)=\chi_r(\pp)g_r(v)$, and $g=W^*\mathrm{diag}(g_1,\dots,g_n)W$ is the required
coboundary.
\end{proof}

\begin{remark}[The series is the primary object]\label{VI:rem:primary}
Without commutation of the Frobenius classes there is no cocycle on $\cR_S$ to be a
coboundary, and the dynamical half of the Kakutani dichotomy has no content; a
non-abelian dynamical model would need a non-abelian acting semigroup, and Chapter~\ref{ch:V} of the
series shows what that costs. What survives, and what everything below uses, is the
Dirichlet series $\Xi_\beta(\rho,S)$ itself, with the following formal properties.
\end{remark}

\begin{lemma}[Additivity and submultiplicativity]\label{VI:lem:additivity}
For representations $\rho,\rho'$ of dimensions $n,n'$:
\begin{align}
\bigl\|I_{n+n'}-(\rho\oplus\rho')(g)\bigr\|_{\HS}^2&=\bigl\|I_n-\rho(g)\bigr\|_{\HS}^2+\bigl\|I_{n'}-\rho'(g)\bigr\|_{\HS}^2, \label{VI:eq:oplus}\\
\bigl\|I_{nn'}-(\rho\otimes\rho')(g)\bigr\|_{\HS}^2&\le2n'\bigl\|I_n-\rho(g)\bigr\|_{\HS}^2+2n\bigl\|I_{n'}-\rho'(g)\bigr\|_{\HS}^2, \label{VI:eq:otimes}\\
\bigl\|I_n-\rho^\vee(g)\bigr\|_{\HS}&=\bigl\|I_n-\rho(g)\bigr\|_{\HS}. \label{VI:eq:dual}
\end{align}
\end{lemma}

\begin{proof}
\eqref{VI:eq:oplus} is block-diagonality of $I-(\rho\oplus\rho')$. For \eqref{VI:eq:otimes} write
$I-A\otimes B=(I-A)\otimes I+A\otimes(I-B)$ and use
$\|X\otimes I_{n'}\|_{\HS}^2=n'\|X\|_{\HS}^2$, $\|A\otimes Y\|_{\HS}^2=n\|Y\|_{\HS}^2$ for
$A$ unitary, together with $\|X+Y\|^2\le2\|X\|^2+2\|Y\|^2$. \eqref{VI:eq:dual} is
$\rho^\vee(g)=\overline{\rho(g)}$ and invariance of the Hilbert--Schmidt norm under entrywise
conjugation.
\end{proof}

\begin{remark}[Numerical check]\label{VI:rem:num1}
Lemma~\ref{VI:lem:trace} and \eqref{VI:eq:otimes} were checked on Haar-random unitaries, $200$ in
each of the dimensions $n=1,2,3,5$ and $500$ pairs with $(n,n')=(2,3)$, with maximal error
$3.6\cdot10^{-15}$ and no violation; a check of algebra, recorded because the constants
of \S\ref{VI:sec:L} depend on it.
\end{remark}

\section{The obstruction is a quotient, not a dual subgroup}\label{VI:sec:tannaka}

\begin{proposition}[Failure of the abelian formulation]\label{VI:prop:failure}
Let $G$ be a non-abelian compact group. Then the set of irreducible representations
$\widehat G$ carries no group structure for which $\rho\mapsto\Xi_\beta(\rho,S)$ is a
homomorphism-like invariant: $\rho_1\otimes\rho_2$ is in general reducible, so
$\widehat G$ is a based set with a fusion rule and not a group, and the annihilator
$\Xi_\beta^\perp$ of a subset of $\widehat G$ is undefined. In particular the statement
``$\mathrm{KMS}_\beta\cong\Prob(G/\Xi_\beta^\perp)$'' has no meaning as written.
\end{proposition}

\begin{proof}
Immediate from the definitions; the abelian case is exactly the case where every
$\rho_1\otimes\rho_2$ is again irreducible.
\end{proof}

\begin{theorem}[Tannakian obstruction]\label{VI:thm:tannaka}
Let $\cT_\beta(S)$ be the full subcategory of $\Rep(G)$ on the objects $\rho$ with
$\Xi_\beta(\rho,S)<\infty$. Then $\cT_\beta(S)$ contains the trivial representation and is
closed under direct sums, subobjects, quotients, tensor products and duals; it is therefore a
Tannakian subcategory of $\Rep(G)$, and
\[
\cT_\beta(S)=\Rep\bigl(G/N_\beta\bigr),
\qquad
N_\beta:=\bigcap_{\rho\in\cT_\beta(S)}\ker\rho ,
\]
a closed normal subgroup of $G$. When $G$ is abelian, $N_\beta=\Xi_\beta^\perp$ and the
description agrees with [Ch.~\ref{ch:Main}, Thm.~3.6].
\end{theorem}

\begin{proof}
The trivial representation gives all summands zero. Closure under $\oplus$, $\otimes$ and
duals is Lemma~\ref{VI:lem:additivity}: for $\otimes$ the bound \eqref{VI:eq:otimes} converts two
convergent series into a convergent one, and $n,n'$ are fixed. For subobjects, if
$\rho=\rho_1\oplus\rho_2$ then \eqref{VI:eq:oplus} gives
$\Xi_\beta(\rho_i,S)\le\Xi_\beta(\rho,S)$, and by unitarity every subobject is a direct
summand; quotients are duals of subobjects. Fullness holds by construction. A full
subcategory of $\Rep(G)$ closed under these operations and containing the trivial object is
the representation category of a quotient of $G$, by Tannaka--Krein for compact groups: the
matrix coefficients of its objects span a $*$-subalgebra of the algebra of representative
functions, closed under the coproduct, hence the representative functions of $G/N$ with $N$
the common kernel. Normality and closedness of $N_\beta$ hold because each $\ker\rho$ is a
closed normal subgroup.

For $G$ abelian, $\cT_\beta$ is generated by the characters in $\Xi_\beta$ and
$\bigcap_{\chi\in\Xi_\beta}\ker\chi=\Xi_\beta^\perp$.
\end{proof}

\begin{remark}[What the theorem does and does not assert]\label{VI:rem:tannakawarn}
Theorem~\ref{VI:thm:tannaka} is a statement about representations and Dirichlet series; it
identifies the natural non-abelian analogue of $\Xi_\beta$ as the quotient $G/N_\beta$. It
does \emph{not} assert that any $C^*$-dynamical system has $\Prob(G/N_\beta)$ as its
$\mathrm{KMS}_\beta$ simplex. That is the content of \S\ref{VI:sec:algebra}, and it is open.

We also restrict to finite-dimensional continuous unitary representations of a compact
group. For a profinite $G$ these are the Artin representations, with finite image. An
$\ell$-adic representation is not unitary and is not covered; the $GL_2$ case of \S5 is
handled through the Sato--Tate group, a compact Lie group, and its Satake classes.
\end{remark}

\section{$L$-functions, with the correct constants}\label{VI:sec:L}

\begin{theorem}[Euler dictionary]\label{VI:thm:euler}
For $\beta>1/2$ and $\rho$ of dimension $n$ unramified outside $S$,
\begin{equation}\label{VI:eq:euler}
\Xi_\beta(\rho,S)=2n\log\zeta_{K,S}(\beta)-2\,\Re\log L_S(\beta,\rho)+O(1),
\end{equation}
where $L_S(s,\rho)=\prod_{\pp\in S}\det\bigl(I-\rho(\Frob_\pp)\Ns\pp^{-s}\bigr)^{-1}$ and the
implied constant depends only on $n$ and $\beta$.
\end{theorem}

\begin{proof}
$\log L_S(s,\rho)=\sum_\pp\sum_{m\ge1}\tfrac1m\Tr\bigl(\rho(\Frob_\pp)^m\bigr)\Ns\pp^{-ms}$,
and the terms with $m\ge2$ are bounded by $n\sum_\pp\sum_{m\ge2}\Ns\pp^{-m\beta}=O_n(1)$ for
$\beta>1/2$. Hence
$\Re\log L_S(\beta,\rho)=\sum_\pp\Ns\pp^{-\beta}\Re\Tr\rho(\Frob_\pp)+O(1)$, and similarly
$\log\zeta_{K,S}(\beta)=\sum_\pp\Ns\pp^{-\beta}+O(1)$. Substituting into \eqref{VI:eq:trace}
gives \eqref{VI:eq:euler}.
\end{proof}

\begin{remark}[Two constants that are easy to get wrong]\label{VI:rem:constants}
The natural guess $\log\zeta_{K,S}(n\beta)-\Re\log L_S(\beta,\rho)=\Xi_\beta(\rho,S)+O(1)$ is
incorrect in two ways. The zeta factor is evaluated at $\beta$, not at $n\beta$: the $n$ in
\eqref{VI:eq:trace} enters as the \emph{coefficient} $2n$ of a single Dirichlet series, coming
from $\Tr I_n=n$, not as a dilation of its argument. And both terms carry a factor $2$, from
$\|I-U\|^2_{\HS}=2(n-\Re\Tr U)$. For $n=1$ the correct identity reduces to
$\Xi_\beta=2\log\zeta_{K,S}(\beta)-2\Re\log L_S(\beta,\chi)+O(1)$, consistent with
[Ch.~\ref{ch:Main}, Rem.~3.11].
\end{remark}

\begin{corollary}[The full prime set]\label{VI:cor:full}
Let $S=P_K$ and $\rho$ irreducible and nontrivial, and suppose $L(s,\rho)$ is holomorphic
and nonzero at $s=1$ --- Hecke--Landau for Artin representations, Jacquet--Shalika
\cite{JS} for the standard $L$-function of a cuspidal automorphic representation. Then
$\Xi_\beta(\rho,P_K)=\infty$ for every $0<\beta\le1$: the obstruction is absent,
$N_\beta=G$, and the would-be simplex is a point.
\end{corollary}

\begin{proof}
For $\beta>1$, \eqref{VI:eq:euler} gives $\Xi_\beta(\rho,P_K)=2n\log\zeta_K(\beta)+O(1)$, since
$\log L(\beta,\rho)$ stays bounded as $\beta\downarrow1$ by the hypothesis on $L(1,\rho)$;
so $\Xi_\beta(\rho,P_K)\to\infty$ as $\beta\downarrow1$. As $\Xi_\beta$ is non-increasing
in $\beta$, being a Dirichlet series with non-negative coefficients, $\Xi_\beta=\infty$ for
all $\beta\le1$.
\end{proof}

The convergent case $\Xi_1(\rho,P_K)<\infty$ would require $\Re\log L(\beta,\rho)$ to
diverge like $n\log\zeta_K(\beta)$, i.e.\ a pole of order $n$ at $s=1$, which happens only
for $\rho$ trivial.

\begin{corollary}[Chebotarev criterion, Artin case]\label{VI:cor:cheb}
Let $\rho$ have finite image and let $D_\rho=\{\pp\in S:\rho(\Frob_\pp)\ne I_n\}$. Then
\[
\Xi_\beta(\rho,S)<\infty\iff\sum_{\pp\in D_\rho}\Ns\pp^{-\beta}<\infty .
\]
\end{corollary}

\begin{proof}
The image is finite, so $\{\|I-\rho(g)\|_{\HS}^2:g\in G,\ \rho(g)\ne I\}$ is a finite set of
strictly positive numbers, bounded above by $4n$ and below by some $c>0$. Hence the summand
in \eqref{VI:eq:Xi} is between $c\Ns\pp^{-\beta}$ and $4n\Ns\pp^{-\beta}$ on $D_\rho$ and zero
off it.
\end{proof}

\begin{remark}
Corollary~\ref{VI:cor:cheb} is the exact analogue of [Ch.~\ref{ch:Main}, Thm.~3.9] and shows that the
Chebotarev machinery of that paper --- in particular the density computations and the block
constructions of [Ch.~\ref{ch:I}, \S5] --- transfers unchanged to Artin representations, with
$D_\rho$ in place of the detecting set of a character. \emph{It fails for infinite image}:
if $\rho(\Frob_\pp)$ can approach $I_n$, the indicator of $D_\rho$ is not comparable to
$\|I-\rho(\Frob_\pp)\|_{\HS}^2$ and only the quadratic form is correct. The $GL_2$ case below
is exactly of this kind.
\end{remark}

\section{$GL_2$: Sato--Tate forbids symmetry breaking, anomalies permit it}

Let $E/\Q$ be a non-CM elliptic curve, $a_p$ the trace of Frobenius, and
$a_p=2\sqrt p\cos\theta_p$ with $\theta_p\in[0,\pi]$. The $\ell$-adic representation of $E$
is not unitary and does not fit Definition~\ref{VI:def:Xi} directly; what does is the
Sato--Tate group $G=SU(2)$ (the non-CM case, by Serre's open image theorem \cite{Serre})
with the Satake classes $\Frob_p=\mathrm{diag}(e^{i\theta_p},e^{-i\theta_p})$ and $\rho$
its standard representation, to which \S\S\ref{VI:sec:kakutani}--\ref{VI:sec:L} apply verbatim.
Then $n=2$ and
\begin{equation}\label{VI:eq:ST}
\bigl\|I_2-\rho(\Frob_p)\bigr\|_{\HS}^2=2\bigl(2-2\cos\theta_p\bigr)=4\bigl(1-\cos\theta_p\bigr)=8\sin^2(\theta_p/2).
\end{equation}

\begin{theorem}[Sato--Tate forbids the obstruction]\label{VI:thm:ST}
Suppose the $\theta_p$, $p\in S$, are equidistributed in $[0,\pi]$ for the Sato--Tate measure
$\mu_{ST}=\tfrac2\pi\sin^2\theta\,d\theta$ in the sense that
$\sum_{p\in S,\ p\le x}p^{-\beta}f(\theta_p)\sim\bigl(\int f\,d\mu_{ST}\bigr)\sum_{p\in S,\ p\le x}p^{-\beta}$
for continuous $f$. Then
\[
\Xi_\beta(\rho,S)\sim4\sum_{p\in S}p^{-\beta},
\]
which diverges whenever the partition function does. Hence $\rho\notin\cT_\beta(S)$ and no
symmetry breaking arises from Sato--Tate distributed classes, at any $\beta$ with
$\zeta_S(\beta)=\infty$.
\end{theorem}

\begin{proof}
By \eqref{VI:eq:ST} the summand is $f(\theta_p)$ with $f(\theta)=4(1-\cos\theta)$, continuous
and bounded on $[0,\pi]$. Since $\int_0^\pi\cos\theta\,\sin^2\theta\,d\theta=0$,
\[
\int_0^\pi f\,d\mu_{ST}=\frac{8}{\pi}\int_0^\pi(1-\cos\theta)\sin^2\theta\,d\theta=4 ,
\]
and the hypothesis applied to $f$ gives the stated asymptotic.
\end{proof}

\begin{remark}[Numerical check]\label{VI:rem:num2}
$\int_0^\pi(2-2\cos\theta)\,d\mu_{ST}=2.0000000000$ by numerical quadrature on $2\times10^6$
nodes, matching the exact value $2$; the factor $2$ of Lemma~\ref{VI:lem:trace} then gives the
constant $4$ in Theorem~\ref{VI:thm:ST}.
\end{remark}

\begin{theorem}[Anomalous sparse sets exist]\label{VI:thm:anomalous}
Assume Sato--Tate for $E$ --- a theorem for non-CM curves over $\Q$ \cite{Taylor} --- in
the form: for every $\delta>0$ the set $\{p:\theta_p\le\delta\}$ has positive natural
density, hence
$\sum_{p\le x,\ \theta_p\le\delta}p^{-1}\to\infty$. Then there is $S\subset P_\Q$ with
\[
\sum_{p\in S}p^{-1}=\infty
\qquad\text{and}\qquad
\Xi_1(\rho,S)<\infty .
\]
\end{theorem}

\begin{proof}
Build $S=\bigsqcup_{j\ge1}S_j$ greedily. Given $S_1,\dots,S_{j-1}$, the hypothesis with
$\delta=2^{-j}$ gives $\sum_{p\le x,\ \theta_p\le2^{-j}}p^{-1}\to\infty$, so a finite set
$S_j$ of primes with $\theta_p\le2^{-j}$, disjoint from the previous blocks, may be chosen
with $\sum_{p\in S_j}p^{-1}\ge1$. Then $\sum_{p\in S}p^{-1}\ge\sum_j1=\infty$. On the other
hand, by \eqref{VI:eq:ST} and $\sin^2(\theta/2)\le\theta^2/4$,
\[
\Xi_1(\rho,S)=\sum_j\sum_{p\in S_j}p^{-1}\cdot8\sin^2(\theta_p/2)
\le\sum_j2\cdot4^{-j}\sum_{p\in S_j}p^{-1} .
\]
Choosing $S_j$ additionally with $\sum_{p\in S_j}p^{-1}\le2$, which is possible since single
terms $p^{-1}$ can be made arbitrarily small, gives $\Xi_1(\rho,S)\le\sum_j4\cdot4^{-j}<\infty$.
\end{proof}

\begin{remark}[What has and has not been shown]\label{VI:rem:anomalouswarn}
Theorem~\ref{VI:thm:anomalous} produces a set on which the analytic obstruction is finite: the
Sato--Tate classes of a non-CM elliptic curve become ``invisible'' along a deliberately
chosen set of primes with anomalously small angle. By Theorem~\ref{VI:thm:tannaka} applied to
$G=SU(2)$ it follows that $\rho\in\cT_1(S)$ and hence $N_1\subseteq\ker\rho=\{1\}$, so the
Tannakian quotient $G/N_1=SU(2)$ is the whole Sato--Tate group, nontrivial and
non-abelian.

What has \emph{not} been shown is that this residual symmetry acts on a simplex of
$\mathrm{KMS}$ states of any algebra, for the reason of \S\ref{VI:sec:algebra}. The phrase
``non-abelian symmetry breaking'' is therefore premature: we have a non-abelian obstruction
group, not yet a broken phase. The construction also assumes Sato--Tate as an
equidistribution input; for non-CM elliptic curves over $\Q$ this is a theorem, but only the
qualitative positive-density statement is used.
\end{remark}

\section{The algebra is not constructed}\label{VI:sec:algebra}

\begin{remark}[Where the $GL_1$ proof stops]\label{VI:rem:gap}
The classification in [Ch.~\ref{ch:Main}] runs: (a) $\mathrm{KMS}_\beta$ states of
$C(Y_S)\rtimes J_S$ correspond bijectively to $\Ns\aaa^{-\beta}$-scaling measures on $Y_S$,
by the conditional expectation onto the diagonal; (b) such measures form a simplex on which
$G_S$ acts; (c) the ergodic decomposition of that action is computed by the Kakutani
criterion. Steps (a) and (b) require the coefficient algebra to be commutative and the
semigroup to be Ore, so that a canonical expectation exists and the crossed product is a
groupoid algebra.

The candidate $\cA^{GL_n}_{K,S}=C(Y^{(n)}_S)\rtimes M$ fails both at present. The
multiplicative structure $M$ is not a semigroup for the obvious choices: $M_n(J_S)$ is not
closed under multiplication with the required cancellation, and $GL_n$ over a semiring of
ideals has no Ore condition available to us. And the Connes--Marcolli $GL_2$ system \cite{ConnesMarcolli2006}, which
does exist, is not of this shape: its symmetry is the action of $GL_2(\mathbb A_f)$ by
endomorphisms of a noncommutative space of two-dimensional lattices, and its arithmetic is
that of modular functions and the modular field, not of a Galois group acting through
Frobenius classes. There is no non-abelian class field theory realized in it, which is why
its $\mathrm{KMS}$ theory does not take the form $\Prob(G/N)$.

The honest statement is therefore: \S\S\ref{VI:sec:kakutani}--\ref{VI:sec:L} determine what the
answer would have to be, and identify the correct algebraic shape of the obstruction, but the
system whose equilibrium states realize $\Prob(G/N_\beta)$ is not available. Two routes seem
worth trying. One is to look for a groupoid model directly: a groupoid whose unit space
carries a $G$-action with the tail relation of [Ch.~\ref{ch:Main}] and whose isotropy realizes
$N_\beta$; this bypasses the semigroup difficulty entirely. The other is to stay with $GL_1$
but to enlarge the coefficient algebra, replacing $C(Y_S)$ by a noncommutative algebra
carrying $\rho$; here the Kakutani analysis of \S\ref{VI:sec:kakutani} applies verbatim, and the
question is whether the resulting crossed product has a tractable expectation.
\end{remark}

\section{Assessment}

\[
\begin{array}{lll}
\text{Statement} & \text{Status} & \text{Reference}\\\hline
\|I-U\|_{\HS}^2=2(n-\Re\Tr U) & \text{true} & \text{Lem.~2.2}\\
\text{Kakutani dichotomy for }U(n)\text{-cocycles} & \text{iff the classes commute} & \text{Thm.~2.3}\\
\Xi_\beta\ \text{a subgroup of }\widehat G & \text{meaningless for non-abelian }G & \text{Prop.~3.1}\\
\Prob(G/\Xi_\beta^\perp) & \text{meaningless as written} & \text{Prop.~3.1}\\
\cT_\beta=\Rep(G/N_\beta) & \textbf{true; the correct replacement} & \text{Thm.~3.2}\\
\log\zeta_{K,S}(n\beta)-\Re\log L_S=\Xi_\beta+O(1) & \text{false} & \text{Rem.~4.2}\\
\Xi_\beta=2n\log\zeta_{K,S}(\beta)-2\Re\log L_S+O(1) & \textbf{true} & \text{Thm.~4.1}\\
\Xi_\beta(\rho,P_K)=\infty,\ \rho\ne1\ \text{irreducible},\ \beta\le1 & \text{true, via }L(1,\rho)\ne0,\infty & \text{Cor.~4.3}\\
\text{Chebotarev criterion via }D_\rho & \text{true for finite image only} & \text{Cor.~4.4}\\
\text{Sato--Tate}\Rightarrow\text{symmetry breaking} & \text{false: forbids it} & \text{Thm.~5.1}\\
\text{anomalous sparse }S\ \text{with}\ \Xi_1<\infty & \text{true, under Sato--Tate} & \text{Thm.~5.3}\\
\cA^{GL_n}_{K,S}\ \text{constructed} & \text{no} & \S\ref{VI:sec:algebra}
\end{array}
\]

The pattern is the one already familiar from Chapter~\ref{ch:III}: the arithmetic
content --- a Dirichlet series in Frobenius data, an $L$-function on $\Re s=1$ ---
generalizes with no difficulty, because none of it used commutativity. What used
commutativity was, first, the dynamical half of the Kakutani step, which needs a cocycle
on a relation generated by commuting shifts, and second the \emph{duality} converting the
series into a statement about a space of states; the substitute for the latter is
Tannakian rather than Pontryagin, a quotient of the group rather than a subgroup of its
dual, and for the former there is none. The missing ingredient is not a theorem but an
object, namely an algebra whose equilibrium states are indexed by $G/N_\beta$.

Three questions. Is $N_\beta$ the intersection of the kernels of \emph{all} representations
in $\cT_\beta$, or is it already cut out by the irreducible ones of bounded dimension? For
$G$ with an open image in $GL_2(\hat\Z)$ this is a question about which quotients can be
realized by anomalous prime sets. Second, Theorem~\ref{VI:thm:anomalous} constructs
$S$ for a fixed $\rho$; can one construct $S$ realizing a \emph{prescribed} quotient
$G/N$, as [Ch.~\ref{ch:I}, Thm.~5.1] does in the abelian case for prescribed
$\bigoplus\Z/2\Z$? Third, and most important, is the groupoid model of Remark~\ref{VI:rem:gap}
available? Without it this paper describes an obstruction group with nothing yet to obstruct.

